% Chapter 1

\section{Relevant Literature} % Main chapter title

\label{Literature} % For referencing the chapter elsewhere, use \ref{Chapter1} 

%----------------------------------------------------------------------------------------

% Define some commands to keep the formatting separated from the content 
\newcommand{\keyword}[1]{\textbf{#1}}
\newcommand{\tabhead}[1]{\textbf{#1}}
\newcommand{\code}[1]{\texttt{#1}}
\newcommand{\file}[1]{\texttt{\bfseries#1}}
\newcommand{\option}[1]{\texttt{\itshape#1}}

%----------------------------------------------------------------------------------------
With this paper we want to argue that agent based models can be a good complement to existing economic modelling practices like equilibrium modelling for policy makers. Agent based models can relax a number of assumptions which have been disputed since the Great Financial Crisis \citep{fagiolo_macroeconomic_2017}. Some of these assumptions are the equilibrium assumption, the representative agent assumption, rational expectations, and linearity or mild non-linearity. By relaxing these assumptions we can perhaps address the inability of macroeconomic models to forecast economic crises.\\

The type of macroeconomic models that is currently dominant in the academic literature and in the central banks are equilibrium type models where dynamics are generated by exogenous shocks. These models are mostly based on the general equilibrium concept of \cite{walras_elements_2013}.  The previous generation of models of Keynesian type endogenous business cycle models, see for instance \cite{kaldor_model_1940} and \cite{hicks_contribution_1950}, fell out of favour as they showed too regular business cycles. \\

Two major findings started the shift from endogenous business cycle models to exogenous shock equilibrium models. The first finding was the proof of the existence of the Walrasian general equilibrium by \cite{arrow_existence_1954} and \cite{mckenzie_equilibrium_1954} independently. The second finding was the work of \cite{slutzky_summation_1937} which showed that a sum of random variables looked very similar to the business cycles that occurred. \cite{hendry_propagation_1995} developed this theory further by showing that a combination of a shock and a propagation mechanism could result in even more realistic business cycles. The combination of these two findings resulted in the next generation of models named real business cycle models or dynamic stochastic general equilibrium models which mostly consists out of two essential ingredients, namely the exogenous shock process and a propagation mechanism. The propagation mechanism in the equilibrium models is the centralized market that is assumed to be in equilibrium by means of a Walrasian auctioneer. \\

However, \cite{sonnenschein_market_1972}, \cite{debreu_economies_1970}, \cite{mantel_characterization_1974} showed that the excess demand function which resulted from equilibrium models with individual preferences could take any shape. This implies that market processes will not necessarily lead to an unique and stable fixed point even when economic agents are optimizing utility \citep{ackerman_still_2002}. Despite this finding, the \cite{lucas_econometric_1976} critique resulted in the dominance of the real business cycle equilibrium models of which one of the first models was developed by \cite{kydland_time_1982}. \\

The problem of non-uniqueness was solved by assuming that the \cite{blanchard_solution_1980} conditions hold\footnote{If these conditions hold, then the system of equations is saddle path stable around the steady state of the model. This assumption results in that after a shock agents instantaneously adjust their behaviour such that the economy jumps to the unique stable path that converges back to the steady state.}. The assumption that the Blanchard-Kahn conditions holds is equivalent to assuming the model to always be locally stable and unique. Dynamics that come out of equilibrium models usually consist out of the following steps: Firstly, the economy is pushed out of the steady state by a shock. Secondly, the economy follows a certain dynamical path which is dictated by the propagation mechanism which is the centralized market in equilibrium. Finally, the model will converge back to the same steady state\footnote{Software packages like Dynare force the Blanchard-Kahn conditions to hold. Otherwise the software will not give you any results.}.\\

The combination of the development of the mathematical techniques, the easiness to use, and the empirical performance of these models resulted in that these models became dominant in the economic profession. Keynesian frictions were added later for which an important example is the \cite{smets_shocks_2007} model of the European economy.\\

\subsection{Disequilibrium}
The use of equilibrium and a Walrasian auctioneer has been disputed before. \cite{patinkin_involuntary_1949} created a model with profit maximization subject to an output constraint where output is lower than the general equilibrium level. Similarly, \cite{clower_keynesian_1969} created a dual decision model where households faced an additional constraint during unemployment, namely that the amount of labour they desired to sell was less than they were able to sell. Consequently, demand thus depended on realized sales of labour instead of desired sales of labour. This resulted in a violation of Walras' law (excess demand has price zero) and consequently the economy was in disequilibrium.  \cite{leijonhufvud_keynesian_1968} used the dual decision model of \cite{clower_keynesian_1969} to built a model without a Walrasian auctioneer. He argued that it was not necessary to introduce price rigidities to have an economy in disequilibrium but simply that prices adjusted slower than quantities or that prices are sufficiently sticky. \cite{barro_general_1971} build on \cite{patinkin_involuntary_1949} and \cite{clower_keynesian_1969} where firms maximized profits subject to an output constraint and households maximized consumption subject to a labour supply constraint. This leads to an economy in general disequilibrium where involuntary unemployment could coexist with "non-excessive" real wages equal to the general equilibrium wage which occurred if both commodity and labour markets were in equilibrium. 

\subsection{Heterogeneity}
Equilibrium models have been developed that depart from the typical representative agent\footnote{Gorman's aggregation theorem stipulates that under quasi-linear utility functions the utility of a population can be represented by a single utility function. This however requires that not a single agent has a different form of utility function. We therefore view the representative agent assumption as a good didactic method but not realistic.}. The first generation of equilibrium models with household heterogeneity are the models that were developed by \cite{bewley_stationary_1987}, \cite{huggett_risk-free_1993}, and \cite{aiyagari_uninsured_1994}. It turns out that computing the equilibrium distribution is quite difficult and that algorithms had to be developed to calculate the dynamics of the models. \cite{krusell_income_1998} did so by using a boundedly rational approach where agents approximate the distribution of wealth by the mean of the distribution. As a result only the mean of the distribution was needed to calculate the equilibrium distribution. Consequently, the common belief of economists was that it was not necessary to develop models that have heterogeneity incorporated. \\

This view changed after the financial crisis where older generation models were found to be insufficient. Heterogeneity was found to be important for policy implication which was supported by new types of distributional equilibrium model such as the HANK model of \cite{kaplan_monetary_2018}. One of the important contributions of this model was the ability to incorporate varying marginal propensities to consume which were found empirically by \cite{fagereng_mpc_2021}. Although these models had a large impact upon policy choices, the models are also criticized of having only mild heterogeneity \citep{fagiolo_macroeconomic_2017}.

Additionally, HANK models are mostly mean-field game models \citep{achdou_income_2022}, which are N-player Nash games where the number of players converge to infinity. Mean field games have one particular important assumptions which allows a N-player Nash game to converge to a continuum player Nash game. This assumption is that each agent has a negligible influence on the other players \citep{cardaliaguet_master_2019}. This implies that situations with spillover effects such as debt cascades\footnote{A good example of a model of debt cascades is the model of \cite{eisenberg_systemic_2001}. Their paper models debt cascades of defaulting banks which form a network with each other where the links are debt claims.} cannot be modelled by mean field games as for these models agents have non-negligible influence on other agents. \footnote{Agent based models find themselves therefore in the middle between mean-field game models and representative agent models. Representative agent models have one agent which determines all the dynamics whereas mean-field games have infinite number of agents where agents have negligible influence. Agent based models can have millions of agents but individual agents can also affect the dynamics of the model.}
\subsection{Expectations}
Most macroeconomic models use the rational expectations hypothesis which was proposed by  \cite{muth_rational_1961}. Following the \cite{lucas_econometric_1976} critique most macroeconomic models switched from adaptive expectations to rational expectations. Most of these models use the expected utility hypothesis from \cite{von_neuman_theory_1953}. However, since then many paradoxes have been found that show that the expected utility hypothesis is an inadequate theory in some cases, such as the \cite{ellsberg_risk_1961} paradox and the \cite{Allais_1953} paradox.\\

\cite{simon_models_1957} proposed an alternative hypothesis known as the bounded rationality hypothesis. Simon argued that people do not have enough processing capacity to make fully rational decisions and suggested that people use heuristics to make decisions. Bounded rationality could both refer to bounded rationality in decision making or bounded rationality in expectation formation. \cite{gabaix_behavioral_2020} showed that a standard New Keynesian model had different policy implications if agents used sparse dynamic programming instead of fully rational dynamic programming. \cite{farhi_monetary_2017} adapt the standard New Keynesian model by letting agents level-k thinking to form expectations. \cite{brock_rational_1997} show that in cobweb type demand-supply model with switching between naive and rational expectations can lead to (chaotic) endogenous fluctuations. \\

The approach of bounded rationality has been critiqued as there are a multitude of approaches which some call the "wilderness of bounded rationality" \citep{sims_macroeconomics_1980}. However, much work was done to empirically validate expectation formation strategies with Learning-to-Forecast experiments. For instance, \cite{anufriev_evolutionary_2012} tested the theory of \cite{brock_rational_1997} using a Learning-to-Forecast experiment and found that in the behaviour of subjects observed in laboratory can be explained very well by an extended version of the theory of \cite{brock_rational_1997}. See \cite{hommes_behavioral_2021} for an overview of these experiments. Additionally, \cite{coibion_how_2018} use survey data to analyse how firms form expectations and find large heterogeneity in firm expectations and inattention to new information which also supports the approach of boundedly rational expectations.

\subsection{Non-linearities and Indeterminacy}
Equilibrium models are often nonlinear when set up mathematically. However, due to computational complexity the model is linearized around the steady state. Several techniques exist that use estimation techniques such as the perturbation method or the projection method but however are similar in nature as they solve the model only locally \citep{den_haan_comparison_2010}. Distributional equilibrium models use similar techniques \citep{marimon_computation_2001}. Moreover, transitions between steady states are computed ad hoc by changing the parameters of the underlying model to force the steady state to move to a different point in the state space \citep{marimon_computation_2001}. By solving the model using linearization, many interesting dynamics are thrown away. For instance, \cite{skiba_optimal_1978} shows with concave-convex production function multiple equilibria can exist. Similarly, \cite{brock_rational_1997} have non-linearities in their model which are derived from heterogeneous expectations whereas \cite{hommes_dynamics_1994} has a nonlinear supply curve with possible dynamics such as k-cycles and chaos instead of convergence to a locally stable and unique steady state.\\

 Under non-linearities multiple equilibria can exist \citep{hommes_behavioral_2013}. Multiple equilibria are not necessarily frowned upon by economists\footnote{The \cite{sonnenschein_market_1972},\cite{mantel_characterization_1974},\cite{debreu_economies_1970} theorem shows that the number of equilibria is finite but not necessarily unique. Only local uniqueness can be guaranteed.}, but the problem lies in selecting an equilibrium. This is called indeterminacy. As mentioned before, equilibrium models usually assume local stability and uniqueness by using stringent restriction such as the \cite{blanchard_solution_1980} conditions \citep{farmer_indeterminacy_2020}. This is done to avoid indeterminacy. However, equilibrium selection is done in arbitrary manner. \cite{arthur_chapter_2006} frames the problem differently and argues that it is not a problem of equilibrium selection but of equilibrium formation. By taking the initial conditions of an economy into account it is possible to see which of the equilibria is selected by simply following the dynamics of the model given the initial conditions. \cite{brock_rational_1997} show for instance that two coexisting period 4 cycles\footnote{A point $a$ is a periodic point of period $n$ or a n-cycle provided $f^n(a)=a$ and $f^j(a) \neq a$ for $0<j<n$. Here, $f^n$ denotes applying the function $f$ $n$ times. Hence, a 4-cycle is a orbit of a system where the system returns to the same point after four evaluations.} can exist. Which 4-cycle occurs is dependent on the initial conditions of the model. Hence, instead of choosing an equilibrium (equilibrium selection) to solve indeterminacy, an alternative approach is to take the initial conditions of an economy and follow the model to which equilibrium it will converge (equilibrium formation).\\
 
Using equilibrium formation also allows endogenous transition from one type of dynamics to another. For instance, \cite{asano_emergent_2021} demonstrate with a heterogeneous agent model that endogenously transitions can occur from a stable equilibrium to an oscillatory regime when a tipping point is crossed over. By opting for equilibrium formation instead of equilibrium selection, the modellers let the model select endogenously which dynamics occur.
\subsection{Agent Based Models}
Typical ingredients of agent based models are heterogeneity, bounded rationality, disequilibrium, and non-linearities. \cite{axtell_agent-based_nodate} argue that agent based models often posses the following features. An agent represents an individual or a group with local information. Furthermore, agents' behaviour are specified by simple rules which are conditional on the state of the agent and global states. The dependency of the agents' behaviour on the agents' states results in heterogeneous behaviour. The agents are embedded in an external environment such as markets or networks which determines the interactions of the agents. Moreover, the model is analyzed by means of simulation which data is consequently analyzed by means of statistical and visual analysis. Aggregate states are not pre-specified but emerge from the interaction of the agents. Finally, size matters in case of agent based models. For instance, a firm that has a size of 100,000 employees behaves differently than a firm with  100 employees. \citep{axtell_agent-based_nodate}. Agent based models could also be seen as the generalization of network models. The stock variables of an agent can be considered as nodes whereas the interaction between the agents could be seen as the edges between the nodes. As agent based models are dynamic, the networks that are generated are dynamic as well which can be useful to analyse dynamics such as systemic risk.\\

Even though equilibrium models exist which include some of these ingredients, to the best of our knowledge there exist none that combine these ingredients. One of the most well known agent based models is the model of segregation of \cite{schelling_models_1969} who showed that a small preference for similar neighbours result in endogenous segregation. Initial economic agent based models focused on microeconomics. For instance, \cite{albin_decentralized_1992} showed with an agent based model and decentralized markets without a Walrasian auctioneer that equilibria could occur which would have larger wealth inequality than a market with a Walrasian auctioneer. \\



Macroeconomic agent based models where initially limited by computing power until high performance computing was more accessible. The EurACE model developed by \cite{deissenberg_eurace_2008} was one of the first models to use high performance computing and featured ten thousand consumers and firms which generated a plethora of macroeconomic phenoma. This model was further developed to allow for policy analysis by \cite{dawid_euraceunibi_2012}.
\cite{riccetti_agent_2015}  developed macroeconomic agent based models further by adding banks. They found that out of their model endogenous dynamics emerge such as business cycles, endogenous growth, unemployment rate fluctuations, leverage cycles, bank defaults, and financial instability. Stability transgresses endogenously to unstable dynamics and vice versa similar to the \cite{minsky_financial_1999} financial instability hypothesis. \cite{assenza_emergent_2015} created an agent based model with capital and credit and consumption and capital firms. Subsequently, \cite{caiani_agent_2016} developed a stock-flow consistent macroeconomic agent based model which introduced credit networks. More macroeconomic agent based models have been developed. See \cite{dawid_agent-based_2018} for an overview.\\

The previous models tried mostly to create internal validity, that is the dynamics generated by the models match the observed data by means of matching moments and correlations. \cite{poledna_economic_2019} innovates by achieving external validity by calibrating an agent based model with Austrian national accounting data taking all economic entities into account and included firms, households, banks, a central bank and a government sector. \cite{poledna_economic_2019} compared the forecasting capabilities of the agent based model with vector auto-regressive models and a dynamic stochastic general equilibrium model and found that the agent based model could compete against the benchmark models. Furthermore, the advantage of using this calibration strategy is that we do not need to worry about equilibrium selection but we simply take the current state of the economy and watch which dynamics occur given the initial state of the economy and thus use equilibrium formation. Another advantage of this calibration technique is that we do not need to use a burn-in period to ensure that the model is stable like other agent based models do such as the models of \cite{caiani_agent_2016} and \cite{riccetti_agent_2015}. This can save a large amount of computation costs. Finally, by the usage of this calibration approach the choice of parameter values is pinned down by empirical evidence instead of endogenous validity.